% The first seeding run written up as an A3 report, in the form the Tutors harness writes for a release.
\begin{tikzpicture}[
    font=\scriptsize,
    panel/.style={draw=jotblue!60, line width=0.6pt, fill=white, text width=7.1cm, align=left,
                  inner sep=5pt, anchor=north west},
    head/.style={font=\scriptsize\bfseries, text=jotblue},
    step/.style={draw=#1, line width=0.7pt, fill=#1!7, rounded corners=2pt, minimum height=0.55cm,
                 minimum width=1.5cm, align=center, inner sep=2pt},
    arr/.style={-{Stealth[length=2mm]}, line width=0.7pt, muted},
  ]
  % sheet
  \draw[jotblue, line width=1pt, fill=wash] (0, 0) rectangle (15.4, -10.6);
  \node[anchor=north west, font=\footnotesize\bfseries, text=jotblue] at (0.25, -0.2)
    {A3: six seeded mistakes passed the gate};
  \node[anchor=north east, font=\scriptsize, text=muted] at (15.15, -0.25)
    {receivers at \texttt{8c4a74a}, 21 September 2026; countermeasures measured 29 September};

  % left column
  \node[panel] (bg) at (0.25, -0.8) {
    \textbf{\textcolor{jotblue}{Background}}\\
    The three receivers enforce structure with architecture rules and a short manifest.
    Seeding twelve canonical mistakes into each and running the single verify command,
    the gate caught 30 of 36.};
  \node[panel] (cur) at (0.25, -2.55) {
    \textbf{\textcolor{jotblue}{Current condition}}\\[3pt]
    \begin{tikzpicture}[baseline]
      \node[step=structure] (w) {write};
      \node[step=behaviour, right=0.5cm of w] (g) {gate};
      \node[step=missed, right=0.5cm of g] (r) {review};
      \draw[arr] (w) -- (g);
      \draw[arr] (g) -- node[above, text=missed, font=\tiny]{V8, V10 pass} (r);
    \end{tikzpicture}\\[2pt]
    V8 (a side effect inlined in the service) and V10 (a business rule placed in the route)
    pass every test in Java, Python and TypeScript and reach a human reviewer. The other
    thirty are stopped at the gate.};
  \node[panel] (goal) at (0.25, -5.0) {
    \textbf{\textcolor{jotblue}{Goal}}\\
    Every seeded mistake stopped at the gate, by the earliest mechanism its language allows,
    with no change to application code.};
  \node[panel] (rca) at (0.25, -6.3) {
    \textbf{\textcolor{jotblue}{Root cause (five whys)}}\\
    1. Why did V8 and V10 pass? Every behavioural test passed.\\
    2. Why? Neither mistake changes what the system does.\\
    3. Why did no rule object? The rules say which layer may import which, not which layer
       may log, notify or decide.\\
    4. Why? No rule asserted placement for side effects or decisions.\\
    5. Why in all three languages? Behaviour cannot see placement, and the pattern had no
       structural rule for it. The gap was in the ladder, not in a receiver.};

  % right column
  \node[panel] (cm) at (7.95, -0.8) {
    \textbf{\textcolor{jotblue}{Countermeasures}}\\
    Two \textsc{Executable Rule}s in each receiver: only the events layer may log or notify;
    routes decide nothing. Each carries a \textsc{Teaching Failure} that names where the
    code belongs. V8 and V10 move from L0 and L2 to L4.\\[2pt]
    While there: a requirement \textsc{Closure Check}, with two new seeds (V13, V14) to
    prove it can fail.};
  \node[panel] (chk) at (7.95, -3.45) {
    \textbf{\textcolor{jotblue}{Check}}\\
    Second run: 42 of 42 caught. V7 now found by an architecture rule rather than a
    behavioural test in Python and TypeScript. Cost: 31 to 89 lines of test code and
    80 to 83 words of manifest per receiver. Application code unchanged.};
  \node[panel] (fu) at (7.95, -5.55) {
    \textbf{\textcolor{jotblue}{Follow-up}}\\
    Repeat the pilot against the current receivers with the request in EARS. Raise the
    harness-erosion use cases from L2 to L3. Adopt the release harness on a second product.};
  \node[panel, fill=jotblue!6] (how) at (7.95, -7.3) {
    \textbf{\textcolor{jotblue}{How this sheet is used}}\\
    In Tutors the harness writes an A3 for every release and its findings inform a confidence
    score. In the receivers the countermeasure became part of the gate. The form is the same;
    a team chooses whether it advises or stops the line.};
\end{tikzpicture}
